% ===== tr-01 块6：第 6 章 pp.9--12（原文页 9--12）=====
% [原文 p.9]
\begin{theorem}{6.3.1}[Hein {[32]}]
对 $\check X$ 上任一 Kähler 形式 $\check\omega'$，若 $\int_{\check X}(\check\omega')^2<\infty$ 且 $[\check\omega']\cdot[C]=0$（其中 $[C]$ 为"坏轮"\footnote{可以看出 $[C]$ 之所以被称为"坏轮"，是因为它似乎是 Ricci 平坦度量存在性的一个障碍。}），则存在 $\check X$ 上完备的 Ricci 平坦度量 $\check\omega$，使得
\begin{enumerate}
\item[(1)] $\check\omega^2=\alpha\Omega\wedge\bar\Omega$，其中 $\alpha$ 充分大。
\item[(2)] $\check\omega-\check\omega'$ 是 $d$-恰当的。
\item[(3)] $\check\omega$ 在无穷远处附近以指数速率衰减到某个半平坦度量。
\end{enumerate}
\end{theorem}
% [Tutor Note: 条件 $[\check\omega']\cdot[C]=0$ 是上同调类 $[\check\omega']$ 与坏轮 $[C]$ 的配对为零（原文印作 $[\check\omega']\cdot[C]$，不可读成逗号分隔的两个类）；这正是剔除坏轮方向的障碍所需的条件。]
% [Tutor Note: 定理 6.3.1(1) 中原文印作 $\Omega\wedge\bar\Omega$（$\Omega$ 无 check），按上下文这里的 $\Omega$ 指 $\check X$ 上的全纯体积形式（下一页半平坦情形记作 $\check\Omega_{sf}$）；正文照原样排。]

下面我们解释定理 6.3.1 中出现的半平坦度量。为综述的需要，我们只考虑 $\check D$ 是 $I_d$-纤维的情形。考虑模型空间
\[
\check X_{\mathrm{mod}}=\Delta^*\times\C/\Lambda(u),\qquad \text{其中 } \Lambda(u)=\Z\oplus\Z\frac{d}{2\pi i}\log u,
\]
其中 $u$ 是去心圆盘 $\Delta^*$ 上的坐标。$\check X_{\mathrm{mod}}$ 通过添加 $u=0$ 处的 $I_d$-纤维而具有自然的部分紧化 $\check Y_{\mathrm{mod}}$。$\Delta$ 上任何只在 $u=0$ 处有一条 $I_d$ 型奇异纤维的椭圆纤维化芽，在选取无穷远附近的局部截面之后，都通过 Abel--Jacobi 映射与 $\check X_{\mathrm{mod}}$ 双全纯。考虑全纯体积形式 $\check\Omega_{sf}=\frac{\kappa(u)}{u}dv\wedge du$，其中 $v$ 是椭圆纤维的坐标，$\kappa(u)$ 是 $u$ 的全纯函数。沿用 Greene--Shapere--Vafa--Yau [28] 的工作，我们考虑半平坦度量
\begin{equation}
\begin{aligned}
\check\omega_{sf,\varepsilon}&:=\sqrt{-1}|\kappa(u)|^2\frac{k|\log|u||}{2\pi\varepsilon}\frac{du\wedge d\bar u}{|u|^2}\\
&\quad+\frac{\sqrt{-1}}{2}\frac{2\pi\varepsilon}{k|\log|u||}\bigl(dv+B(u,v)du\bigr)\wedge\overline{\bigl(dv+B(u,v)du\bigr)},
\end{aligned}
\tag{6.3.2}
\end{equation}
其中 $B(u,v)=-\frac{\operatorname{Im}(v)}{\sqrt{-1}u|\log|u||}$，而 $\varepsilon$ 是椭圆纤维的体积。容易看出 $\check\omega_{sf}$ 限制到每条纤维上是平坦度量，并且可以直接验证 $2\check\omega_{sf}^2=\check\Omega_{sf}\wedge\bar{\check\Omega}_{sf}$，即 $\check\omega_{sf}$ 是超 Kähler 度量。
% [Tutor Note: 式 (6.3.2) 的第一项是底方向（$u$ 方向，含 $du\wedge d\bar u/|u|^2$）的度量，第二项是纤维方向（$v$ 方向）的度量；系数中 $|\log|u||$ 是双重绝对值（外层为模、内层为 $|u|$），不要抄成 $\log|u|$。]
% [Tutor Note: $B(u,v)$ 的分母是 $\sqrt{-1}u|\log|u||$（含绝对值），与式 (6.3.2) 的系数一致；$\varepsilon$ 是椭圆纤维的体积，不是参数 $\varepsilon$ 的某种极限。]
% [Tutor Note: $2\check\omega_{sf}^2=\check\Omega_{sf}\wedge\bar{\check\Omega}_{sf}$ 是曲面上超 Kähler（Calabi--Yau）度量的归一化条件：它等价于 $\check\omega_{sf}$ 的 $(2,0)$-部分的模平方由全纯体积形式给出。]

同调群 $H_2(\check X_{\mathrm{mod}},\Z)$ 秩为二，由纤维与"坏轮" $[C]\in H_2(\check X_{\mathrm{mod}},\Z)$ 生成；后者是唯一一个使它在 $H_2(\check Y_{\mathrm{mod}},\Z)$ 中的推前同调于零的同调类。显式地，$[C]$ 可以由曲面 $\{|u|=\text{常数},\operatorname{Im}(x)=\text{常数}\}$ 实现。直接计算表明，这些是相对于 $(\check\omega_{sf},\check\Omega)$ 的特殊拉格朗日环面，并且这些闭链显然定义了 $\check X_{\mathrm{mod}}$ 的一个纤维化。
% [Tutor Note: 坏轮的显式实现 $\{|u|=\text{常数},\ \operatorname{Im}(x)=\text{常数}\}$ 中用的是 $\operatorname{Im}(x)$（$x$ 为纤维方向坐标），不是 $\operatorname{Im}(u)$；它与前面 $v$ 坐标的关系是 $v$ 的虚部，按原文照抄。]

\begin{remark}{6.3.2}
$\check X$ 上的半平坦度量由 $\check X_{\mathrm{mod}}$ 上的半平坦度量 $\omega_{sf,\varepsilon}$ 诱导，同时还需要选取 $\check X$ 在无穷远处附近的一个全纯截面 $\sigma$；这一选取通过把 $\check X$ 的无穷远邻域与 $X_{\mathrm{mod}}$ 等同起来而进入。关键在于，不同的截面选取可能给出截然不同的 Kähler 度量；例如它们可以不是 de Rham 上同调的，也可以不是一致等价的。因此，$\check X$ 上的半平坦度量应当记作 $\omega_{sf,\sigma,\varepsilon}$，以记录对 $\sigma$ 的依赖。然而，为本综述的目的，我们将略去对 $\sigma$ 的依赖。
\end{remark}
% [Tutor Note: 原文此处把局部模型空间印作 $X_{\mathrm{mod}}$（无 check），按上下文应为 $\check X_{\mathrm{mod}}$；正文照原样排。]

% [原文 p.10]
借助 Hein 的稳健估计，与定理 6.2.1 和定理 6.2.3 中相同的论证证明了：上述特殊拉格朗日环面模型的 LMCF 收敛到相对于 $(\check\omega,\check\Omega)$ 的真正特殊拉格朗日，并且可以形变为特殊拉格朗日。总结起来，我们有

\begin{theorem}{6.3.3}
（见 [11, 定理 3.4]）给定有理椭圆曲面 $\check X$ 与一条 $I_d$-纤维 $\check D$，对 Hein 所构造的任一 Ricci 平坦 Kähler 度量，$\check X=\check Y\setminus\check D$ 上都存在特殊拉格朗日纤维化。此外，$\check X$ 的合适超 Kähler 旋转 $\check X'$ 可以通过添加一条 $I_d$-纤维而紧化为有理椭圆曲面 $\check Y'$。
\end{theorem}

\begin{remark}{6.3.4}
若 $\check D$ 是 $I_d$ 以外任何类型的奇异纤维，则由 Mayer--Vietoris 序列，$\check X$ 的无穷远邻域的第二同调群（模去挠元）由椭圆纤维生成。特别地，这些情形下 $\check X$ 不能容许特殊拉格朗日纤维化。
\end{remark}

\section{非标准半平坦度量}
Hein 观察到，由定理 6.3.1 构造的 $\check X$ 上的 Ricci 平坦度量具有与二维 Tian--Yau 空间相同的总体几何量，包括曲率与单射半径的衰减以及体积增长；见式 (6.2.1)。因此自然要问两个问题：
\begin{enumerate}
\item[(Q1)] 作为复流形，$\check X$ 是否是 $X$ 的一个合适的超 Kähler 旋转？
\item[(Q2)] 若 (Q1) 的答案是肯定的，则在超 Kähler 旋转之后，Tian--Yau 度量是否与 Hein 的度量一致？
\end{enumerate}
问题 (Q1) 已在定理 6.2.3 中得到肯定回答，但 (Q2) 的答案却是否定的。确实，假设第二个问题的答案是肯定的。取次数为 $d$ 的弱 del Pezzo 曲面 $Y$ 与一般的光滑反典范除子 $D$，$X=Y\setminus D$。由定理 6.2.3，$X$ 中存在特殊拉格朗日纤维化。此外，$X$ 有一个合适的超 Kähler 旋转 $\check X$，它可以通过添加一条 $I_d$-纤维而紧化为有理椭圆曲面 $\check Y$。若 $\check X$ 上的超 Kähler 度量是 Hein 的度量之一，则由定理 6.3.3，$\check X$ 容许特殊拉格朗日纤维化。作反方向的超 Kähler 旋转，就得到 $X$ 中一个新的特殊拉格朗日纤维化，它不同于我们出发时的那个，且相恰好相差 $\pi/2$。这意味着 $D$ 的基本域可以取成一个矩形，而由于我们一开始取的是 $D$ 的一般选取，这是荒谬的。总之，弱 del Pezzo 曲面与有理椭圆曲面上的特殊拉格朗日纤维化似乎探测到了 $\check X$ 上一些此前未知的超 Kähler 度量。特别地，这表明 $\check X_{\mathrm{mod}}$ 上的超 Kähler 度量比式 (6.3.2) 形式的那些更多。这促使人们寻找 $\check X_{\mathrm{mod}}$ 上 Ricci 平坦度量的新拟设。在 [11, 第 2 节] 中发现，对每个 $b_0\in\R$，存在一个非标准半平坦度量，由下式给出
% [Tutor Note: (Q2) 的反证用的是"相分析"：若 Hein 度量与 Tian--Yau 度量在旋转后一致，则 $X$ 中会出现两个相相差 $\pi/2$ 的特殊拉格朗日纤维化，从而把一般椭圆曲线 $D$ 的基本域压成矩形------这对一般 $D$ 不可能。这是原文"absurd"论证的可验证内容。]

% [原文 p.11]
\begin{equation*}
\begin{aligned}
\check\omega_{sf,b_0,\varepsilon}&:=\sqrt{-1}\frac{|\kappa(u)|^2}{\varepsilon}W^{-1}\frac{du\wedge d\bar u}{|u|^2}\\
&\quad+\frac{\sqrt{-1}}{2}W\varepsilon\bigl(dv+\widetilde{\Gamma}(v,u,b_0)du\bigr)\wedge\overline{\bigl(dv+\widetilde{\Gamma}(u,v,b_0)du\bigr)},
\end{aligned}
\end{equation*}
其中 $W=\frac{2\pi}{k|\log|u||}$，而 $\widetilde{\Gamma}(u,v,b_0)=B(u,v)+\frac{b_0}{2\pi^2}\frac{|\log|u||}{u}$。
% [Tutor Note: 上式第一个括号内 $\widetilde{\Gamma}$ 的变量次序原文印作 $(v,u,b_0)$，第二个括号内印作 $(u,v,b_0)$（且第二个括号整体加横线表示复共轭）；按上下文两者应同为 $(u,v,b_0)$，正文照原样保留。]
% [Tutor Note: 系数 $\frac{2\pi}{k|\log|u||}$ 中的 $k$ 与式 (6.3.2) 中的 $k$ 是同一记号；$|\log|u||$ 同样是双重绝对值。]

\begin{remark}{6.4.1}
若 $2b_0/d\in\Z$，则 $\check\omega_{sf,b_0,\varepsilon}$ 是标准半平坦度量经截面 $\sigma:\Delta^*\to\check X_{\mathrm{mod}}$ 作逐纤维平移之后的结果。否则，$\check\omega_{sf,b_0,\varepsilon}$ 给出 $\check X_{\mathrm{mod}}$ 上一个新的 Ricci 平坦度量。直接计算表明，$\check\omega_{sf,b_0,\varepsilon}$ 限制到椭圆纤维上也是平坦的。换言之，$\check X_{\mathrm{mod}}$ 上存在一族由 $b_0$ 参数化的单参数半平坦度量族，其中离散的一组选取给出 Greene--Shapere--Vafa--Yau 的标准半平坦度量。
\end{remark}

几何上，非标准半平坦度量可以通过沿某个多值截面（其取值甚至可能不可数）给出的逐纤维平移映射，从标准半平坦度量拉回而得到。考虑如下形式的多值截面 $\sigma:\Delta^*\to\check X_{\mathrm{mod}}$：
\[
\begin{gathered}
\sigma:\Delta^*\longrightarrow\check X_{\mathrm{mod}}\\
u\mapsto\Bigl(u,\,h(u)+\frac{c_0}{2\pi i}\log u+\frac{b_0}{(2\pi i)^2}(\log u)^2\Bigr),
\end{gathered}
\]
其中 $b_0,c_0\in\R$，$h(u)$ 是 $\Delta^*$ 上的全纯函数。设 $T_\sigma^*$ 为沿多值截面 $\sigma$ 的逐纤维平移。直接计算表明 $T_\sigma^*\check\omega_{sf}=\check\omega_{sf,b_0,sf}$ 在 $\check X_{\mathrm{mod}}$ 上良定义。值得注意的是，$2b_0/d\in\Z$ 当且仅当 $\sigma$ 延拓为 $\Delta\to\check Y_{\mathrm{mod}}$ 上的截面，见 [20, 引理 3.28]。
% [Tutor Note: 多值截面的第二个分量是 $\frac{c_0}{2\pi i}\log u+\frac{b_0}{(2\pi i)^2}(\log u)^2$，注意分母分别是 $(2\pi i)$ 与 $(2\pi i)^2$，不是 $2\pi$ 或 $2\pi^2$。]
% [Tutor Note: 原文此处"$T_\sigma^*\check\omega_{sf}=\check\omega_{sf,b_0,sf}$"中的下标印作 $b_0,sf$；按定义应为 $\check\omega_{sf,b_0,\varepsilon}$（同族第三个参数是 $\varepsilon$），这是原文的笔误，正文照原样排。]

非标准半平坦度量 $\check\omega_{sf,b_0,\varepsilon}$ 可以作为 $\check D$ 附近的模型度量，用来沿 Hein [32] 的论证在 $\check X$ 上构造完备超 Kähler 度量。记 $[F]\in H_2(\check X,\Z)$ 为纤维的同调类。通过在 $\check D$ 附近选取局部截面，我们可以把 $\check X$ 的无穷远处附近的椭圆纤维化与局部模型 $\check X_{\mathrm{mod}}$ 等同起来。

\begin{theorem}{6.4.2}
（见 [11, 定理 2.16]）设 $\check\omega_0$ 是 $X$ 上的 Kähler 形式，使得
\[
[\check\omega_0]_{dR}\cdot[F]=\varepsilon,\qquad [\check\omega_0]_{dR}\cdot[C]=\frac{2b_0}{\varepsilon},
\]
其中 $b_0\in\R$。则存在 $\check X$ 上的 Kähler 度量 $\check\omega_\alpha$，使得
\[
\check\omega_\alpha=\alpha T_h^*\check\omega_{sf,b_0,\frac{\varepsilon}{\alpha}}
\]
在 $\check D$ 附近成立，并且在 $\check D$ 的某个邻域之外有 $\check\omega_\alpha=\check\omega_0$。这里 $h$ 是 $\Delta^*$ 上的全纯截面，它只依赖于 Bott--Chern 上同调类
% [Tutor Note: 定理条件中 $\check\omega_0$ 的所在空间原文印作 $X$（下一行的 $H^{1,1}_{BC}(X,\R)$ 同），但 $\check\omega_0$ 是 $\check X$ 的 Kähler 形式、且 $H^{1,1}_{BC}$ 依赖于复结构，故按上下文应为 $\check X$；正文照原样排。]

% [原文 p.12]
$[\check\omega_0]_{BC}\in H^{1,1}_{BC}(X,\R)$。此外，对充分大的 $\alpha$\footnote{该条件可以看作：存在性在充分接近大体积极限时成立；见第 6.6.1 节的讨论。}，存在 $\check X$ 上的超 Kähler 度量
\[
\check\omega=\check\omega_\alpha+i\partial\bar\partial\varphi,
\]
使得
\begin{enumerate}
\item[(1)] $\check\omega^2=\alpha^2\check\Omega\wedge\bar{\check\Omega}$；且
\item[(2)] $\check\omega$ 渐近于（可能非标准的）半平坦度量 $T_h^*\check\omega$，其含义是存在常数 $\delta>0$，使得对每个 $k\in\N$ 有
\[
|\nabla^k\varphi|\sim O(e^{-\delta r^{2/3}}),
\]
其中 $r$ 是到 $\check X$ 中某个固定点的距离函数。
\end{enumerate}
\end{theorem}
% [Tutor Note: (1) 中 $\check\omega^2=\alpha^2\check\Omega\wedge\bar{\check\Omega}$ 的 $\alpha^2$ 与定理 6.4.2 中 $\check\omega_\alpha$ 参数的 $\alpha$ 一致；这里的 $\check\Omega$ 带 check（与 p.9 定理 6.3.1(1) 中无 check 的 $\Omega$ 不同）。]
% [Tutor Note: (2) 中原文印作"半平坦度量 $T_h^*\check\omega$"（$\check\omega$ 无下标）；按上下文（与本节开头及定理末尾的 $\alpha T_h^*\check\omega_{sf,b_0,\varepsilon/\alpha}$ 对照），所指应是某个半平坦度量 $T_h^*\check\omega_{sf,b_0,\varepsilon}$。正文照原样排。]

特别地，$\check\omega$ 与 $\check\omega_{sf,b_0,\varepsilon}$ 具有相同的曲率衰减 $|\nabla^k Rm|\le r^{-2-k}$（对每个 $k\in\N$）以及单射半径衰减 $\operatorname{inj}\sim r(x)^{-1/3}$。
% [Tutor Note: 这里的曲率衰减指数是 $r^{-2-k}$（不是式 (6.2.1) 中的 $-\frac{k+2}{n+1}$）：因为本定理处理的是复二维（曲面）模型，故指数退化为 $-2-k$。]

\begin{remark}{6.4.3}
可以证明，对任意 $\varepsilon,b_0$，都存在满足定理 6.4.2 假设的环境 Kähler 形式 $\check\omega_0$；见 [11] 或下文定理 6.6.14 的证明。
\end{remark}

对每个满足 $2b_0/d\in\Q$ 的 $b_0$，存在同调类 $C_{b_0}\in H_2(\check X_{\mathrm{mod}},\Z)$，使得 $\int_{C_{b_0}}\check\omega_{sf,b_0,\varepsilon}=0$，并且 $\check X_{\mathrm{mod}}$ 上存在相对于 $(\check\omega_{sf,b_0,\varepsilon},\check\Omega_{sf})$、以 $C_{b_0}$ 为纤维类的特殊拉格朗日纤维化。与定理 6.3.3 类似，我们有如下定理。

\begin{theorem}{6.4.4}
沿用上述记号。若 $2b_0/d\in\Q$，则 $\check X$ 上存在以 $C_{b_0}$ 为纤维类、相对于定理 6.4.2 中 Ricci 平坦度量的特殊拉格朗日纤维化。此外，$\check X$ 的合适超 Kähler 旋转 $\check X'$ 可以通过添加一条 $I_d$-纤维而紧化为有理椭圆曲面 $\check Y'$。
\end{theorem}

借助非标准半平坦度量，我们现在可以解释 Calabi ansatz 的超 Kähler 旋转。

\begin{theorem}{6.4.5}
（见 [11, 附录 A]）设 $D\cong\C/\Z\oplus\Z\tau$ 是椭圆曲线，其中 $\tau$ 位于上半平面。沿用第 6.2.1 节的记号，考虑 $X_{\C}$ 的超 Kähler 旋转，其 Kähler 形式为 $\check\omega_{\C}$、全纯体积形式为 $\check\Omega_{\C}$，使得对应于 $1\in\Z\oplus\Z\tau$ 的 ansatz 特殊拉格朗日的相为零。则适当选取坐标后有
\[
\check\omega_{\C}=\alpha\check\omega_{sf,b_0,\varepsilon},\qquad \check\Omega_{\C}=\alpha\check\Omega_{sf},
\]
其中 $b_0=-\frac{1}{2}Re(\tau)d$，$\varepsilon=\frac{2\sqrt{2\pi}}{Im(\tau)}$，$\alpha=\sqrt{d\pi Im(\tau)}$。特别地，$\tau\leftrightarrow(b_0,\varepsilon)$ 之间存在双射，即每个（可能非标准的）半平坦度量都可以实现为某个 Calabi ansatz 的某个超 Kähler 旋转，至多相差一个伸缩。
\end{theorem}
% [Tutor Note: 定理 6.4.5 中三个公式不要互相混淆：$b_0=-\frac{1}{2}Re(\tau)d$、$\varepsilon=\frac{2\sqrt{2\pi}}{Im(\tau)}$、$\alpha=\sqrt{d\pi Im(\tau)}$；注意 $\varepsilon$ 的分母是 $Im(\tau)$ 而不是 $Re(\tau)$；此处原文的 Re、Im 排作数学斜体（与 p.6 的正体 $\operatorname{Re},\operatorname{Im}$ 不同），正文照原样。]
% [Tutor Note: "ansatz special Lagrangian corresponding to $1\in\Z\oplus\Z\tau$" 指由椭圆曲线 $D\cong\C/(\Z\oplus\Z\tau)$ 上平移截面给出的 ansatz 特殊拉格朗日；这里 $\tau$ 在上半平面，故 $Im(\tau)>0$。]
